\section{The conditional main theorem}\label{sec:conditional-lift}

The comparison results of Section~\ref{sec:frozen-frontier} are purely order-theoretic. To speak about effective extension operators, such as those that arise from arithmetic, the state space needs an effective structure. The missing canonicalization step then has to be stated as an explicit assumption.

\subsection{Effective setting and contract}

From now on $X$ carries a preorder $\le$ and an effective encoding of its elements by natural numbers, in the sense of mathlib's \texttt{Primcodable} class. An operator $e$ is \emph{computable} if it is computable with respect to this encoding, and \emph{monotone} if $x\le y$ implies $e(x)\le e(y)$. In the intended application $X$ is a set of sentences, the encoding is a standard G\"odel numbering, and $\le$ is provable implication.

Fix a ledger $L$, a domain $D$, a cone $C$, a canonical family $\mathcal{K}$, and an operator $e$ under study.

\begin{definition}[Contract]\label{def:contract}
The \emph{contract} for $(L,D,C,\mathcal{K},e)$ consists of three clauses.
\begin{enumerate}
\item[\textup{(C1)}] \emph{Membership:} $\mathcal{K}\subseteq D$.
\item[\textup{(C2)}] \emph{Admissibility:} every $c\in\mathcal{K}$ is computable and monotone.
\item[\textup{(C3)}] \emph{Canonicalization:} if $e$ is computable, monotone and efficient in $D$ under $L$, then $e$ is aligned with $\mathcal{K}$ on $C$.
\end{enumerate}
\end{definition}

Clauses (C1)--(C3) are hypotheses of the main theorem, not results. They are not proved here for any arithmetic choice of $(X,D,C,\mathcal{K})$. Section~\ref{sec:arithmetic-instance} establishes the corresponding facts in one concrete arithmetic case, up to provable equivalence.

\subsection{Statement}

\begin{theorem}[Conditional canonical representative]\label{thm:conditional-arithmetic-lift}
Let $L$ be invariant on $C$, and let $e$ be computable, monotone and efficient in $D$ under $L$. If the contract \textup{(C1)--(C3)} holds for $(L,D,C,\mathcal{K},e)$, then there is $c\in\mathcal{K}$ such that
\begin{enumerate}
\item $c$ is computable and monotone,
\item $c$ is efficient in $D$ under $L$, and
\item $e\agree{C}c$.
\end{enumerate}
\end{theorem}

\begin{proof}
Clause (C3) makes $e$ aligned with $\mathcal{K}$ on $C$. Lemma~\ref{lem:restricted-alignment}, which uses (C1) and the invariance of~$L$, gives $c\in\mathcal{K}$ that is efficient in $D$ with $e\agree{C}c$. Clause (C2) gives computability and monotonicity of~$c$.
\leanref{\lid{conditional\_arithmetic\_canonicality\_lift}, exported as \PaperMainTheorem}
\end{proof}

The proof is short because the substantive step, canonicalization, is assumed. The value of the theorem lies in its exact hypotheses. Each clause is needed, and the conclusion keeps the link $e\agree{C}c$ to the original operator. The package results of Section~\ref{sec:fixed-package} are not used. They concern growth, while this theorem concerns comparison.

\subsection{Why admissibility must be assumed}

It might seem that $c$ should inherit computability and monotonicity from $e$, since the two agree on the cone. They do not. Agreement on $C$ says nothing about values outside $C$, while computability and monotonicity are global properties.

\begin{proposition}[Clause (C2) cannot be dropped]\label{prop:admissibility-necessary}
Take $X=\N$ with its usual order and encoding, $C=\{0\}$, $D=\mathcal{E}$, the constant ledger $Y\equiv K\equiv 0$, and $e=\mathrm{id}$. Let $b(n)=2$ for odd $n$ and $b(n)=0$ for even $n$, and put $\mathcal{K}=\{b\}$. Then $L$ is invariant on $C$, the identity is computable, monotone and efficient, and (C1) and (C3) hold. Yet no $c\in\mathcal{K}$ is monotone.
\end{proposition}

\begin{proof}
Under a constant ledger no operator dominates another, so every member of $D$ is efficient. (C1) holds because $D$ contains every operator. (C3) holds because $b(0)=0=\mathrm{id}(0)$. The map $b$ is computable but not monotone, since $1\le 2$ while $b(1)=2>0=b(2)$.
\leanref{\lid{canonical\_admissibility\_is\_necessary} in \texttt{tests/lean/MathReview.lean}}
\end{proof}

The regression file also shows that invariance of the ledger cannot be dropped from Theorem~\ref{thm:efficiency-transfer}. It further shows that an efficient operator need not have any representative in a given family, for instance the empty one. Theorem~\ref{thm:ledger-boundary} states the general form of the second point.

\subsection{The intended arithmetic reading}

The intended reading takes $X$ to be arithmetic sentences, the order to be provable implication over a base theory, $e$ to be a recursive monotone extension operator, and $\mathcal{K}$ to be a family of canonical extensions by consistency or reflection statements, in the tradition of~\cite{Turing1939,Feferman1962,BeklemishevReflection,PakhomovWalsh2021}. Two cautions apply.

First, literal agreement $e(x)=c(x)$ on sentence codes is too strong for arithmetic. Two operators that produce provably equivalent but syntactically different sentences do not agree in this sense. The natural notion is agreement up to provable equivalence. Passing to the Lindenbaum quotient would turn equivalence into equality, but no computable normal form for that quotient is available, so computability on codes would be lost. Section~\ref{sec:arithmetic-instance} therefore works with equivalence explicitly. It proves the analogue of Theorem~\ref{thm:efficiency-transfer} for ledgers invariant under provable equivalence on the cone.

Second, canonicalization for \emph{all} recursive monotone operators cannot be expected. Walsh has shown that comparison with the iterates of the consistency operator breaks down for general recursive monotone operators. There are such operators, even with $\Pi_1$ values, whose strength oscillates cofinally between different iterates~\cite{Walsh2022Evitable}. Clause (C3) must therefore be established for restricted classes. Section~\ref{sec:arithmetic-instance} does this for one explicit operator and for one class defined by Walsh's hypotheses.
